% basic nastaveni
\documentclass[12pt]{article}
\setlength{\parindent}{0pt}
\setlength{\parskip}{1em}
\usepackage[utf8]{inputenc}
\usepackage[czech]{babel}
\usepackage[T1]{fontenc}
\usepackage{graphics}
\usepackage{caption}
\usepackage{hyperref}

% matematicke package
\usepackage{amsmath}
\usepackage{amsfonts}
\usepackage{amssymb}
\usepackage{geometry}
\geometry{margin=2.5cm}
\usepackage{pgfplots}
\pgfplotsset{compat=1.18}

\title{Maturitní otázka číslo 14: Zobrazení v rovině}
\author{}
\date{}

\begin{document}
\maketitle

\subsection*{Co to je zobrazení}
Zobrazení $f$ v rovině je předpis, který každému bodu $X$ roviny přiřazuje právě jeden bod $X'$ této
roviny. Bod $X$ se nazývá vzor a bod $X'$ obraz. Zapisujeme to jako $f(X) = X'$ nebo $f : X.
\mapsto X'$. Zobrazení se nazývá prosté právě tehdy, když každé dva různé body $X, Y$ zobrazí na dva
různé body $X', Y'$ roviny. To je důležité pokud bychom chtěli dělat zobrazení inverzní, pokud by
zobrazení nebylo prosté, nešlo by zjistit původní vzor.

Máme-li zobrazení $f$ v rovině, samodružný bod je takový bod $X$, který se zobrazí sám na sebe, tedy
platí $f(X) = X$. Samodružná přímka $p$ je taková přímka, která se zobrazí do sebe sama, tedy platí
$f(p) = p$.

\subsection*{Shodná zobrazení}
Zobrazení se nazývá shodné zobrazení, pokud pro každé dva body $X, Y$ roviny a jejich obrazy $X', Y'$
platí $|XY| = |X'Y'|$. Obrazem každého geometrického útvaru je geometrický útvar schodný, tedy obraz
se dá tak přemístit, že se vzor s ním kryje. Specifickým typem shodného zobrazení je identické
zobrazení, které každému bodu $X$ roviny přiřazuje ten stejný bod $f(X) = X$.

Schodná zobrazení se dále dělí na přímou a nepřímou schodnost. Představme si trojúhelník, které ma
svoje vrcholy popsané proti směru hodinových ručiček. Příma schodnot nechá směr pojmenování
trojúhelníku stejný a patří mezi ní identita, posunutí, otočení a středová souměrnost. Nepřímá
schodnot vytvoří trojúhelník který bude schodný, ale bude pojmenován po směru hodinových ručiček. Jde
si to představit jako kdyby ho nekdo vzal a otočil jako palačinku. Mezi nepřímou schodnost patří
osová souměrnost.
\subsubsection*{Středová souměrnost}
Středová souměrnost $S(S)$ se středem v bodě $S$ je zobrazení ve kterém se bod $S$ zobrazí na bod
$S' = S$ a každý bod $X$, který není $S$ se zobrazí tak, že bod $S$ je středem úsečky $XX'$. Platí
tedy, že $|XS| = |SX'|$. Bod $S$ se nazývá střed souměrnosti. Středová souměrnost se středem
souměrnosti v bodě $S$ se zapisuje jako $S(S) : X \mapsto X'$ a udává že bod $X'$ je obrazem bodu $X$.

Samodružným bodem středové souměrnosti je střed souměrnosti a samodružné přímky jsou všechny
přímky, které prochází středem souměrnosti.

Při konstrukcích se často využívá to, že jsou nějaké útvary středově souměrné. Patří mezi ně
například čtverec, obdélník, kružnice, kosočtverec nebo kosodélník. Středem souměrnosti kružnice je
její střed. U ostatních vyjmenovaných útvarů je to průsečík jejich úhlopříček.
\subsubsection*{Posunutí}
Posunutí $T(AB)$ neboli translace určené orientovanou úsečkou $AB$ je zobrazení, ve kterém se zobrazí
bod $X$ na bod $X'$ tak, že orientované úsečky $AB$ a $XX'$ jsou rovnoběžné a mají stejnou délku
a směr. Orientovaná úsečka se také nazývá vektor posunutí. Posunutí vektorem $\vec{AB}$ se zapisuje
jako $T(AB) : X \mapsto X'$ a udává že bod $X'$ je obrazem bodu $X$. 

Samodružný bod existuje pokud je vektor posunutí roven nule a budou jím všechny body roviny. Jinak
posunutí samodružný bod nemá. Samodružná přímka vzniká pokud je přímka rovnoběžná s vektorem posunutí.
\subsubsection*{Otočení}
Otočení $R(S, \alpha)$ neboli rotace určená bodem S a orientovaným úhlem $\alpha$ je zobrazení, ve
kterém se bod $S$ zobrazí na bod $S' = S$ a každý bod $X$ který není $S$ na bod $X'$ tak, že
$|XS| = |X'S|$ a orientované úhly $XSX'$ a $\alpha$ mají stejnou velikost a jsou stejně orientované.
Bod $S$ se nazývá střed otočení a orientovaný úhel $\alplha$ se nazývá úhel otočení. Otočení se
středem otočení v bodě $S$ a orientovaným úhlem $\alpha$ se zapisuje jako $R(S, \alpha) : X \mapsto
X'$ a udává že bod $X'$ je obrazem bodu $X$.

Středová souměrnost je speciálním typem otočení a to o $180^\circ$.

Soudružný bod existuje v bodu otáčení a nebo je tvoří celá rovina, pokud je $\alpha = k \cdot
360^\circ, k \in \mathbb{Z}$. Soudružnou přímku získáme pokud pokud je znovu $\alpha = k \cdot
360^\circ, k \in \mathbb{Z}$ a nebo pokud přímka prochází středem otáčení a $\alpha = 180^\circ + k
\cdot 360^\circ, k \in \mathbb{Z}$.
\subsubsection*{Osová souměrnost}
Osová souměrnost $O(o)$ s osou $o$ je zobrazení, ve kterém se zobrazí každý bod $X$ na ose $o$ na bod
$X' = X$ a každý bod $X$ který neleží na ose na bod $X'$ tak, že úsečka $XX'$ je kolmá na osu a střed
$S$ úsečky $XX'$ leží na ose. Platí tedy, že $|XS| = |SX'|$. Přímka o se nazývá osa souměrnosti.
Osovou souměrnost s osou $o$ se zapisuje jako $O(o) : X \mapsto X'$ a udává že bod $X'$ je obrazem
bodu $X$.

Všechny samodružné body leží na ose souměrnosti. Samodružné přímky jsou všechny přímky kolmé na osu
souměrnosti a samotná osa.

Při konstrukcích se často využívá to, že jsou nějaké útvary osově souměrné. Patří mezi ně
například čtverec, obdélník, kružnice, kosočtverec nebo kosodélník.

\subsection*{Podobná zobrazení}
to do
zmin ze kazda shodnost je podobne zobrazeni s koeficientem 1, ne naopak

\end{document}
